% Bilaga B, satt ur info/quartus_workflow.md.

\chapter{Quartus Prime Lite och DE0-CV}
\label{app:quartus}

Det här är bokens verktygsreferens för syntes och hårdvara. Genom \chapref{c1:ch} till \ref{c18:ch}
är den \textbf{lärarvänd}: Quartus- och DE0-CV-arbetet demonstreras från katedern, med början i
bromsassistenten i \chapref{c1:ch}, och ingen deltagare behöver Quartus eller ett eget kort.

\begin{aside}
\Chapref{c19:ch} är undantaget, och det är det som räknas. Bring-upen demonstreras först och lämnas
sedan över: varje grupp syntetiserar och programmerar sin egen \code{can_spi_node} och demonstrerar
sin egen tvånodslänk. Läs därför \secref{app:quartus:project} till \ref{app:quartus:program} som
deltagarvända \textemdash{} du behöver \secref{app:quartus:project} och \ref{app:quartus:sources} för
att skapa projektet och lägga till källfilerna, och \secref{app:quartus:compile} är kompileringen som
producerar den \file{.sof} \secref{app:quartus:program} programmerar \textemdash{} och läs
\secref{app:quartus:linux} \emph{innan} du kopplar in något, eftersom det är avsnittet som avgör om
Programmer över huvud taget kan se kortet.
\end{aside}

Ingenting som \textbf{bedöms} kräver hårdvaran (bilaga~\ref{app:project}). Bring-upen bedöms på
simuleringen bakom den, inte på om ledningen fungerade just den dagen. Allt som beskrivs här slutar
på riktig FPGA-hårdvara, verifierad för hand. Det är ett annat lager än övningarna, som du verifierar
på din egen laptop med en självkontrollerande testbänk under GHDL (bilaga~\ref{app:sim}).

% ------------------------------------------------------------------------------------------
\section{Vad som behövs}
\label{app:quartus:need}

\begin{itemize}
  \item \textbf{Quartus Prime Lite} (gratisutgåvan): tillverkarens FPGA-verktyg, som täcker syntes,
    placering och routning (”fitting”), och generering av programmeringsfil.
  \item \textbf{Enhetsstöd för Cyclone V}: ett separat tilläggspaket. DE0-CV:s FPGA,
    \code{5CEBA4F23C7N}, är en Cyclone V-krets, och Quartus installerar inte varje familjs stöd som
    förval.
  \item \textbf{USB-Blaster-drivrutin}: DE0-CV programmeras över USB med USB-Blaster-protokollet.
    Drivrutinen installeras oftast tillsammans med Quartus, men behöver ibland en manuell knuff
    (\secref{app:quartus:trouble}).
  \item Ett \textbf{Terasic DE0-CV}-kort och en USB-kabel.
\end{itemize}

Quartus Prime Lite finns för både Windows och Linux, och konstruktionsstegen nedan är desamma
oavsett värdoperativsystem. \textbf{Att få kortet synligt för Programmer är det inte}, och det är där
varje värdspecifikt problem i den här kursen bor: \secref{app:quartus:linux} täcker Linux och WSL,
\secref{app:quartus:trouble} täcker Windows.

\begin{warning}
En rekommendation direkt, eftersom den sparar en kväll: kursens GHDL-verktygskedja kör på WSL eller
Ubuntu, men \textbf{WSL2 kan inte se en USB-enhet alls utan extra arbete}
(\secref{app:quartus:linux}). Är du på Windows, installera Quartus nativt på Windows och behåll WSL
för GHDL. De två behöver aldrig se varandra; de delar bara \file{.vhd}-filerna, som ligger i
Windows-filsystemet hur som helst.
\end{warning}

% ------------------------------------------------------------------------------------------
\section{Att installera Quartus Prime Lite}
\label{app:quartus:install}

\begin{enumerate}
  \item Hämta \textbf{Quartus Prime Lite Edition} från tillverkarens nedladdningsportal. Sök efter
    ”Quartus Prime Lite download”: Intels FPGA-verksamhet blev \textbf{Altera} igen 2024, så sidorna
    har flyttat och äldre djuplänkar ruttnar. Det som är stabilt att söka på är produktnamnet, inte
    URL:en.
    \begin{itemize}
      \item Välj \textbf{Lite}-utgåvan, som är den fria och inte behöver någon licensfil. Standard
        och Pro täcker inte Cyclone V gratis.
      \item Vilken någorlunda färsk Lite-utgåva som helst fungerar; kursen beror inte på en viss
        delversion.
      \item Välj Windows- eller Linux-installeraren efter var du tänker köra den, och läs varningen
        ovan om du är på Windows med WSL.
    \end{itemize}
  \item Se under installationen till att \textbf{Cyclone V} är valt bland enhetsfamiljerna. Är
    Quartus redan installerat utan det, lägg till Cyclone V-stödet efteråt, antingen genom att köra
    installeraren igen eller via Quartus egen ”Additional Software”-installerare.
  \item När DE0-CV-kortet ansluts över USB första gången ska Windows upptäcka USB-Blastern och
    antingen installera drivrutinen automatiskt eller be dig peka ut Quartus egna drivrutinsfiler
    (vanligen under Quartus installationskatalog). Hittar Windows ingen drivrutin automatiskt, se
    \secref{app:quartus:trouble}.
\end{enumerate}

\subsection{Två fallgropar i Windows 11}
\label{app:quartus:win11}

Ett par av Windows 11:s säkerhetsförval kan störa installationen, eller programmeringen av kortet
senare:

\begin{itemize}
  \item \textbf{Om installeraren inte ens öppnar sig:} leta upp \code{Smart App Control} i Windows
    Säkerhet (sök efter det i startmenyn), stäng av det och starta om.
  \item \textbf{Om allt installeras snyggt men kortet ändå inte går att programmera}
    (\secref{app:quartus:program}): gå till \code{Windows Security → Device security → Core
    isolation} och slå av \textbf{Memory integrity}.
\end{itemize}

\lecfig[0.78]{info/images/mem_integrity.png}{\code{Memory integrity} under Windows 11:s
inställningar för kärnisolering.}{app:quartus:fig:memint}

Ingendera inställningen rör Quartus i sig; båda är allmänna härdningsförval i Windows 11 som råkar
vara strikta nog att störa USB-åtkomst på drivrutinsnivå.

% ------------------------------------------------------------------------------------------
\section{Att skapa ett nytt projekt}
\label{app:quartus:project}

\begin{enumerate}
  \item Öppna Quartus och starta \textbf{New Project Wizard} (\code{File → New Project Wizard}).
  \item Välj en arbetskatalog och ett projektnamn. Quartus använder projektnamnet som förvalt namn
    på toppnivåentiteten också, så låt det matcha din VHDL-entitets namn för att slippa skriva över
    det senare (för \file{or_gate.vhd}: döp projektet till \code{or_gate}).
  \item Välj \textbf{Empty Project} som projekttyp. Kursen använder varken Quartus IP-katalog eller
    någon projektmall.
  \item Hoppa över att lägga till filer i det här steget (\secref{app:quartus:sources} täcker det),
    om du inte redan vet exakt vilken fil det gäller; endera går bra.
  \item På sidan för enhetsval, välj exakt den krets som sitter på DE0-CV: filtrera på familjen
    \textbf{Cyclone V}, och leta sedan upp och välj enheten \code{5CEBA4F23C7N}.
  \item Hoppa över sidan för EDA-verktyg (lämna den på ”none”); kursen kör sina simuleringar direkt i
    GHDL i stället för genom Quartus.
  \item Avsluta guiden.
\end{enumerate}

% ------------------------------------------------------------------------------------------
\section{Att lägga till sin VHDL}
\label{app:quartus:sources}

\begin{itemize}
  \item Lägg till källfilen: \code{Project → Add/Remove Files in Project...}, bläddra fram din
    \file{.vhd}-fil och lägg till den.
  \item Har konstruktionen mer än en fil (en toppnivåentitet plus en eller flera delkomponenter),
    lägg till dem allihop. Quartus löser upp instansieringshierarkin automatiskt så länge varje fil
    finns i projektet.
  \item Bekräfta att toppnivåentiteten är rätt satt: \code{Project → Set as Top-Level Entity}, med
    din huvudfil markerad. För en enfilskonstruktion sker det automatiskt.
\end{itemize}

% ------------------------------------------------------------------------------------------
\section{Att tilldela pinnar i Pin Planner}
\label{app:quartus:pins}

Varje port i din entitet behöver en fysisk pinne på FPGA:n innan Quartus kan producera en fungerande
programmeringsfil; annars finns det ingen väg från konstruktionens portar till en verklig
strömbrytare eller lysdiod på kortet.

\begin{enumerate}
  \item Kör \code{Processing → Start → Start Analysis \& Elaboration} först, och öppna sedan Pin
    Planner via \code{Assignments → Pin Planner}. I ett alldeles nytt projekt är portlistan tom
    tills Quartus har elaborerat konstruktionen och känner till dess portar.
  \item Var och en av entitetens portar dyker upp som en rad. Fyll i kolumnen \textbf{Location} med
    pinnbeteckningen för den strömbrytare eller lysdiod du vill använda.
  \item \textbf{En strömbrytare eller en lysdiod är fritt val. En klocka är det inte.} För
    \code{SW}, \code{KEY} och \code{LEDR} går det bra att ta vilka som helst som är bekväma:
    konstruktionen bryr sig inte om vilken strömbrytare som driver vilken ingång, bara att varje
    port når en. \code{CLOCK_50} är en bestämd pad på kretsen och går inte att flytta. På DE0-CV är
    den \textbf{\code{PIN_M9}}, oscillatorn på 50 MHz. Varje konstruktion från \chapref{c3:ch} och
    framåt är klockad, så från det kapitlet är det här det enda pinnummer som är värt att kunna utan
    att slå upp. Bekräfta det mot ditt eget korts filer (steg 4) innan du förlitar dig på det.
  \item \textbf{Skriv inte av pinntabellen för hand.} Tio \code{LEDR}-pinnar avskrivna ur en PDF är
    tio chanser att kasta om två siffror, och felet det ger är det otäckaste i
    \secref{app:quartus:trouble}: en ren kompilering som beter sig obegripligt.
    \begin{itemize}
      \item DE0-CV:s system-CD innehåller en \file{DE0_CV.qsf} där varje kortpinne redan är namngiven
        och tilldelad.
      \item \code{Assignments → Import Assignments…}, peka ut den filen, och varje
        \code{CLOCK_50}-, \code{SW[n]}-, \code{KEY[n]}-, \code{LEDR[n]}- och
        \code{GPIO_0[n]}-tilldelning landar på en gång.
      \item Döp din entitets portar så att de matchar de signalnamnen så återstår ingenting att
        tilldela. Det är därför \chapref{c19:ch}:s bring-up döper sina portar till \code{CLOCK_50},
        \code{KEY}, \code{SW}, \code{LEDR} och \code{GPIO_0} i stället för till något snyggare.
      \item Pin Planner (och kortets pinntabeller i användarmanualen) är fortfarande vägen att
        kontrollera en enskild pinne, eller tilldela en som \file{.qsf}-filen inte täcker.
    \end{itemize}
  \item Sätt kolumnen \textbf{I/O Standard} om Quartus inte härleder ett vettigt förval för pinnen du
    valde. DE0-CV:s allmänna I/O är 3,3 V LVTTL/LVCMOS. Att importera kortets \file{.qsf} (steg 4)
    sätter det här också, vilket är ännu ett skäl att föredra den vägen.
  \item Spara tilldelningarna (de skrivs in i projektets \file{.qsf}-fil). Varje ändring här kräver
    omkompilering (\secref{app:quartus:compile}) innan den får effekt; Quartus lappar inte i
    efterhand en \file{.sof} som kompilerades före ändringen.
\end{enumerate}

\subsection{En klocka behöver ett tidskrav också}
\label{app:quartus:sdc}

Att tilldela \code{CLOCK_50} en pinne talar om för Quartus \emph{var} klockan kommer in. Det säger
inte hur fort den går, och tills du säger det har tidsanalysen ingenting att kontrollera mot och
rapporterar en obegränsad konstruktion. Lägg till en \file{.sdc}-fil i projektet med en rad:

\begin{console}
create_clock -name CLOCK_50 -period 20.000 [get_ports CLOCK_50]
\end{console}

20 ns är 50 MHz. Med den på plats blir tidsanalysen i \code{Processing → Start Compilation}
meningsfull och rapporten ger dig ett verkligt Fmax, vilket är siffran \secref{c5:sec:fmax}
skickar dig hit för att läsa. Utan den saknas Fmax, eller är inte att lita på.

% ------------------------------------------------------------------------------------------
\section{Att kompilera konstruktionen}
\label{app:quartus:compile}

\code{Processing → Start Compilation} (eller ▶ i verktygsraden) kör hela flödet: analys och syntes,
placering och routning, tidsanalys, och hopsättning till en programmeringsfil.

Läs kompileringsrapporten när den är klar, i synnerhet \textbf{Warnings}: en oansluten port, en
ofullständig känslighetslista, eller en latch som härletts där du menade kombinatorik dyker alla upp
där som varningar snarare än fel. De är lätta att missa om man bara kontrollerar att bocken blev
grön.

En lyckad kompilering producerar en \file{.sof}-fil (SRAM Object File) i projektets
\file{output_files}-katalog. Det är den \secref{app:quartus:program} laddar ner till kortet. Den är
flyktig: den konfigurerar FPGA:ns SRAM direkt och försvinner när strömmen bryts. Kursen har inget
steg som bränner konfigurationen permanent.

% ------------------------------------------------------------------------------------------
\section{Att programmera DE0-CV}
\label{app:quartus:program}

\begin{enumerate}
  \item Anslut DE0-CV till datorn över USB och slå på det.
  \item Öppna Programmer: \code{Tools → Programmer}.
  \item Välj USB-Blastern under \textbf{Hardware Setup}. Den ska hittas automatiskt om drivrutinen
    installerades rätt (annars \secref{app:quartus:trouble}).
  \item Se till att \file{.sof}-filen från \secref{app:quartus:compile} står i listan och att dess
    kryssruta \textbf{Program/Configure} är ibockad. Står den inte där, använd \code{Auto Detect}
    eller lägg till filen för hand.
  \item Klicka \textbf{Start}. Förloppsindikatorn ska nå 100 \% och rapportera att det gick bra.
  \item Testa konstruktionen direkt på kortet: växla de strömbrytare du tilldelade i
    \secref{app:quartus:pins} och bekräfta att lysdioden följer entitetens sanningstabell.
\end{enumerate}

% ------------------------------------------------------------------------------------------
\section{Felsökning}
\label{app:quartus:trouble}

\begin{description}[style=unboxed,leftmargin=0pt,itemsep=0.8ex]
  \item[USB-Blastern syns inte i Programmer.] Kontrollera Enhetshanteraren i Windows efter en okänd
    enhet eller en med drivrutinsfel, och peka manuellt ut USB-Blaster-drivrutinen under Quartus
    installationskatalog.
  \item[Kortet går inte att programmera, utan något tydligt fel.] Slå av \code{Memory integrity}
    under \code{Core isolation} på Windows 11 (\secref{app:quartus:win11}). Det är den enskilt
    vanligaste orsaken när drivrutinen väl är installerad.
  \item[Installeraren startar inte.] Kontrollera \code{Smart App Control} under Windows Säkerhet
    (\secref{app:quartus:win11}).
  \item[Kompileringen går igenom men konstruktionen beter sig fel på kortet.] Kontrollera
    tilldelningarna i Pin Planner (\secref{app:quartus:pins}) \emph{innan} du misstänker din VHDL.
    En omkastad eller saknad pinntilldelning ger ett kort som kompilerar rent och beter sig
    obegripligt, eftersom det inte är kopplat som du tror.
  \item[Fitter rapporterar otilldelade pinnar som fel.] Vissa projektinställningar behandlar det som
    fatalt snarare än som en varning. Tilldela antingen varje port i Pin Planner, eller, för en
    kvarglömd oanvänd port tidigt i utvecklingen, säg åt Quartus hur resten ska behandlas:
    \code{Assignments → Device → Device and Pin Options… → Unused Pins}. För DE0-CV är
    \textbf{”As input tri-stated (with weak pull-up)”} det säkra värdet. Förvalet i vissa flöden
    driver oanvända pinnar låga, vilket på ett kort där pinnar når kontaktlister och kringkretsar är
    ett sätt att slåss med en annan drivare. Lämna inte portar tyst otilldelade och förvänta dig att
    Quartus gissar.
\end{description}

% ------------------------------------------------------------------------------------------
\section{Att programmera från Linux och från WSL}
\label{app:quartus:linux}

\Secref{app:quartus:trouble} är Windows. På Linux är drivrutinen ingen drivrutin alls: Programmer
talar med USB-Blastern genom libusb, och behöver rättighet att göra det.

\subsection*{Linux, nativt}
Ur lådan visar Programmer ingen hårdvara, eftersom enhetsnoden ägs av root. Ge den en udev-regel med
raderna nedan i \file{/etc/udev/rules.d/51-usbblaster.rules}:

\begin{console}
# Altera USB-Blaster
SUBSYSTEM=="usb", ATTR{idVendor}=="09fb", ATTR{idProduct}=="6001", MODE="0666"
SUBSYSTEM=="usb", ATTR{idVendor}=="09fb", ATTR{idProduct}=="6002", MODE="0666"
SUBSYSTEM=="usb", ATTR{idVendor}=="09fb", ATTR{idProduct}=="6003", MODE="0666"
# USB-Blaster II
SUBSYSTEM=="usb", ATTR{idVendor}=="09fb", ATTR{idProduct}=="6010", MODE="0666"
SUBSYSTEM=="usb", ATTR{idVendor}=="09fb", ATTR{idProduct}=="6810", MODE="0666"
\end{console}

Ladda om reglerna och dra sedan ur och i kortet igen:

\begin{shell}
sudo udevadm control --reload-rules && sudo udevadm trigger
\end{shell}

Quartus cachar sin hårdvarulista i en bakgrundsdemon, så om Programmer fortfarande inte visar något,
starta om den:

\begin{shell}
killall jtagd 2> /dev/null
jtagconfig            # ska nu lista USB-Blastern och kedjan av enheter
\end{shell}

\code{jtagconfig} är den snabbaste vägen att skilja ett rättighetsproblem från ett kabelproblem:
listar den kortet är hårdvaran i sin ordning och problemet sitter inuti Quartus; gör den det inte,
stanna och lös det innan du rör Programmer.

\subsection*{WSL2}
Det här är fallet kursens egen verktygskedja går rakt in i, så det är värt att säga rent ut:
\textbf{WSL2 har ingen USB-genomkoppling.} Kortet är osynligt för en WSL2-installation hur många
udev-regler du än skriver, eftersom kärnan aldrig ser enheten. Det finns ingen inställning att ändra.
Två vägar ut:

\begin{itemize}
  \item \textbf{Kör Quartus nativt på Windows} (rekommenderas). Behåll WSL för GHDL, som inte behöver
    någon hårdvara alls. Dina \file{.vhd}-filer ligger i Windows-filsystemet och båda verktygen läser
    dem. Det är vägen kursen förutsätter, och den med minst rörliga delar.
  \item \textbf{Koppla vidare enheten med \code{usbipd-win}}, om du specifikt vill ha Quartus inuti
    WSL. Från en PowerShell som administratör på Windows-sidan:

\begin{console}
winget install usbipd
usbipd list                      # hitta USB-Blasterns BUSID
usbipd bind   --busid <BUSID>    # en gång per enhet, består
usbipd attach --wsl --busid <BUSID>
\end{console}

    Enheten dyker då upp inuti WSL och udev-regeln ovan gäller. Den måste kopplas om efter varje
    urdragning och varje omstart av WSL, och Quartus grafiska gränssnitt behöver dessutom WSLg eller
    en X-server. Görbart, men tre saker till som kan gå fel på en redovisningsdag.
\end{itemize}
